%==============================================================================
% Appendix D: Computational notes and normalization chain
%==============================================================================
\section{Computational Notes and Normalisation Chain}
\label{app:D}

This appendix collects the normalisation chain and the descriptor separation
used in the body.  Only the minimum auxiliary material needed to keep the
main dictionary distinct from the secondary readings is recorded.

\subsection{Geometric Response and Descriptors}
\label{app:D:descriptor}

We distinguish the geometric response itself from the descriptor that
expresses it in another language.  The primary object of the body is always
the geometric response dictionary, while the KK reading and the
Matsubara-style thermal reading are secondary interpretations attached to
it.

\begin{itemize}
  \item \textbf{geometric response}: the primary object of the body.
  \item \textbf{reduced descriptor}: a description in lower-dimensional
        language.
  \item \textbf{thermal descriptor}: a description in thermal language.
\end{itemize}

This distinction prevents conflating the response dictionary with the
re-reading languages.  For example, $\KMxw$ is simultaneously a native
geometric quantity, a reduced kinetic baseline, and a thermal stiffness
scale; this does not imply that the three are mutually identical.
Likewise, $\etaAPS$ and $\Ntop$ both admit lower-dimensional / thermal
descriptors, but as observable families they are distinct.

\subsection{Normalisation Chain}
\label{app:D:chain}

The normalisation patch of the CZ boundary term should be read along the
chain
\begin{equation}
  \mathrm{NY} = \dd \Ctor
  \;\Longrightarrow\;
  \int_{X_4}\mathrm{NY} = \int_{M_3}\Ctor
  \;\Longrightarrow\;
  \Ntop := \frac{1}{4\pi^2}\int_{M_3}\Ctor
  \;\Longrightarrow\;
  \SCZ[M_3] = \frac{k_{\mathrm{cl}}}{4\pi^2}\int_{M_3}\Ctor.
\end{equation}
The torsional boundary term, the frame-bundle-normalised torsional charge,
and the inflow action are thus seen to be different representations within
the same $\Ctor$ sector.  The fact that the $1/(4\pi^2)$ normalisation
corresponds to $T_R = 1$ (frame-bundle Chern-class unit) and that
$\Ntop = 6 r_0^2$ in $\Sthree$ is derived in Appendix~\ref{app:E:E9}.

The APS spectral entry is not part of this chain.  $\etaAPS$ belongs to the
spectral family carried by $\Csig$ and must be treated as an observable
distinct from $\Ntop$ and $\SCZ$, which are defined directly from $\Ctor$.
The reason for adopting the three-family taxonomy is precisely that the
family distinction is not erased by this normalisation chain.

\subsection{Reduced-Side Consistency Statements}
\label{app:D:consistency}

Under the KK reduction reading, the lower-dimensional parity-odd language
admits the following three consistency statements:
\begin{equation}
  k_{3D}^{\mathrm{tor\text{-}CS}} = \frac{1}{2}\,k_{4D}^{\NY},
  \qquad
  Q_{\mathrm{best}} := \frac{\eta(0)}{\Delta k_{\mathrm{matter}}} = 1,
  \qquad
  k_q\in\mathbb{Z}.
\end{equation}
The first is the correspondence between the 4D Nieh--Yan coefficient and
the 3D torsion-Chern--Simons coefficient (KK reduction normalisation
identity; Appendix~\ref{app:E:E11}; verification:
\texttt{proofs/kk\_normalization.py}).  The second is a normalisation-free
quotient using a spectral quantity and a matter shift (matter-minimal
inflow audit; Appendix~\ref{app:E:E12}).  The third is the integrality
condition of the reduced Chern--Simons level (formal quantisation on the
frame bundle; Appendix~\ref{app:E:E13}).

In this paper these are not added to the boundary observable family as new
entries.  Instead, they are positioned as consistency statements that arise
when reading the Euclidean geometric response dictionary already given in
the body in lower-dimensional language.  We do not claim a complete
identification between these reduced-side statements and the native
geometric quantities.
